\chapter{Electrocardiogram (ECG/EKG)}

\section*{What Is This Test?}
An electrocardiogram records the heart's electrical activity through electrodes placed on the skin. It shows heart rate and rhythm and can provide evidence of conduction abnormalities, ischemia, prior myocardial injury, chamber enlargement, electrolyte effects, or medication effects.

\section*{Alternative Names}
ECG; EKG; 12-Lead ECG; Electrocardiography; Resting ECG

\section*{Why It Is Ordered}
Evaluation of chest pain, palpitations, fainting, shortness of breath, dizziness, suspected arrhythmia, or possible heart attack; assessment of known heart disease; and baseline or medication-related monitoring when clinically indicated.

\section*{How This Test Is Performed}
The patient lies still while adhesive electrodes are placed on the chest, arms, and legs. The machine records several electrical views for a few seconds. The test is painless and usually takes only a few minutes. Longer monitoring may be performed with a Holter or event monitor when symptoms are intermittent.

\section*{Preparation, Risks, and Limitations}
No fasting is required. Avoid applying lotion or oil to the skin immediately beforehand, and tell the technician about implanted devices or skin sensitivity. The test has no electrical risk; removing electrodes may briefly irritate the skin. A normal resting ECG does not exclude intermittent arrhythmia or all forms of coronary disease.

\section*{References}
\begin{enumerate}
  \item MedlinePlus. Electrocardiogram.
  \item American Heart Association. Electrocardiogram (ECG or EKG).
\end{enumerate}

\clearpage
